%%%%%%%%%%%%%%%%%%%%%%%%%%%%%%%%%%%%%%%%%%%%%%%%%%%%%%%%%%%%%%%%%%%%%%%%
% Shared preamble for both build roots (main-lncs.tex, main-preprint.tex).
%
% Contains everything that is identical across the two layouts: math and
% graphics packages, colours, the \val{} generated-values macro, and the
% single-source hooks (\tightvspace, \abskeywords) that each root may
% override. Format-specific setup (document class, one/two column, title
% block, bibliography style) lives in the individual roots.
%
% This file is \input after the document class (and, for the preprint
% build, after \usepackage{neutral2col}), so package load order is:
%   class / style  ->  natbib (preprint only)  ->  these packages
%   ->  hyperref  ->  cleveref.
%%%%%%%%%%%%%%%%%%%%%%%%%%%%%%%%%%%%%%%%%%%%%%%%%%%%%%%%%%%%%%%%%%%%%%%%

\usepackage[T1]{fontenc}
\usepackage{microtype}
\usepackage{graphicx}
\usepackage{booktabs}
\usepackage{amsmath}
\usepackage{amssymb}
\usepackage{mathtools}
\usepackage{xcolor}
\usepackage{tikz}
\usepackage{float}        % provides the [H] float placement used in figures
\usepackage{subcaption}   % provides appendix subfigure panels and \captionsetup

\DeclareMathOperator*{\argmax}{arg\,max}

%%%%%%%%%%%%%%%%%%%%%%%%%%%%%%%%
% Colours (used by figures/drift_types.tex and the \ph mark-up macro)
%%%%%%%%%%%%%%%%%%%%%%%%%%%%%%%%
\definecolor{light_orange}{RGB}{255,229,191}
\definecolor{light_red}{RGB}{253,195,183}
\definecolor{light_gray}{RGB}{198,230,233}
\definecolor{light_blue}{RGB}{230, 230, 255}
\definecolor{accent_orange}{RGB}{255, 128, 2}
\definecolor{accent_blue}{RGB}{178, 178, 255}
\definecolor{accent_blue_dark}{RGB}{102, 102, 255}
\definecolor{light_grey}{gray}{0.95}
\definecolor{phviolet}{RGB}{124,58,237}
\newcommand{\ph}[1]{\textcolor{phviolet}{#1}}

% Generated result values, cited through \val{<dataset>/<path>}.
% Generated by `drift analysis figures`; do not edit by hand.
% Read a value with \val{<dataset>/<path>}, e.g. \val{arxiv/model/tx_m/decay}.
% An unknown key stops the build rather than printing nothing.
\providecommand{\val}[1]{%
  \ifcsname valstore/#1\endcsname\csname valstore/#1\endcsname
  \else\errmessage{unknown value key: #1}\fi}
% ---- Yearbook ----
\expandafter\def\csname valstore/yearbook/models\endcsname{21}
\expandafter\def\csname valstore/yearbook/cutoffs\endcsname{104}
\expandafter\def\csname valstore/yearbook/seeds\endcsname{5}
\expandafter\def\csname valstore/yearbook/cutoff-first\endcsname{1905}
\expandafter\def\csname valstore/yearbook/cutoff-last\endcsname{2013}
\expandafter\def\csname valstore/yearbook/scratch-models\endcsname{12}
\expandafter\def\csname valstore/yearbook/frozen-models\endcsname{9}
\expandafter\def\csname valstore/yearbook/model/clip_b32_frozen/indist\endcsname{73.2}
\expandafter\def\csname valstore/yearbook/model/clip_b32_frozen/future\endcsname{64.1}
\expandafter\def\csname valstore/yearbook/model/clip_b32_frozen/decay\endcsname{9.1}
\expandafter\def\csname valstore/yearbook/model/cnn_l/indist\endcsname{92.8}
\expandafter\def\csname valstore/yearbook/model/cnn_l/future\endcsname{79.0}
\expandafter\def\csname valstore/yearbook/model/cnn_l/decay\endcsname{13.7}
\expandafter\def\csname valstore/yearbook/model/cnn_m/indist\endcsname{92.3}
\expandafter\def\csname valstore/yearbook/model/cnn_m/future\endcsname{78.4}
\expandafter\def\csname valstore/yearbook/model/cnn_m/decay\endcsname{13.9}
\expandafter\def\csname valstore/yearbook/model/cnn_s/indist\endcsname{87.0}
\expandafter\def\csname valstore/yearbook/model/cnn_s/future\endcsname{74.7}
\expandafter\def\csname valstore/yearbook/model/cnn_s/decay\endcsname{12.3}
\expandafter\def\csname valstore/yearbook/model/convnext_s_frozen/indist\endcsname{79.9}
\expandafter\def\csname valstore/yearbook/model/convnext_s_frozen/future\endcsname{70.6}
\expandafter\def\csname valstore/yearbook/model/convnext_s_frozen/decay\endcsname{9.2}
\expandafter\def\csname valstore/yearbook/model/dinov2_s_frozen/indist\endcsname{72.5}
\expandafter\def\csname valstore/yearbook/model/dinov2_s_frozen/future\endcsname{62.7}
\expandafter\def\csname valstore/yearbook/model/dinov2_s_frozen/decay\endcsname{9.8}
\expandafter\def\csname valstore/yearbook/model/dinov3_s_frozen/indist\endcsname{65.8}
\expandafter\def\csname valstore/yearbook/model/dinov3_s_frozen/future\endcsname{57.9}
\expandafter\def\csname valstore/yearbook/model/dinov3_s_frozen/decay\endcsname{7.9}
\expandafter\def\csname valstore/yearbook/model/eva02_b_frozen/indist\endcsname{72.7}
\expandafter\def\csname valstore/yearbook/model/eva02_b_frozen/future\endcsname{63.4}
\expandafter\def\csname valstore/yearbook/model/eva02_b_frozen/decay\endcsname{9.3}
\expandafter\def\csname valstore/yearbook/model/mae_b_frozen/indist\endcsname{67.4}
\expandafter\def\csname valstore/yearbook/model/mae_b_frozen/future\endcsname{59.0}
\expandafter\def\csname valstore/yearbook/model/mae_b_frozen/decay\endcsname{8.4}
\expandafter\def\csname valstore/yearbook/model/mlp_l/indist\endcsname{84.9}
\expandafter\def\csname valstore/yearbook/model/mlp_l/future\endcsname{72.7}
\expandafter\def\csname valstore/yearbook/model/mlp_l/decay\endcsname{12.2}
\expandafter\def\csname valstore/yearbook/model/mlp_m/indist\endcsname{85.7}
\expandafter\def\csname valstore/yearbook/model/mlp_m/future\endcsname{73.0}
\expandafter\def\csname valstore/yearbook/model/mlp_m/decay\endcsname{12.7}
\expandafter\def\csname valstore/yearbook/model/mlp_s/indist\endcsname{69.3}
\expandafter\def\csname valstore/yearbook/model/mlp_s/future\endcsname{62.3}
\expandafter\def\csname valstore/yearbook/model/mlp_s/decay\endcsname{7.1}
\expandafter\def\csname valstore/yearbook/model/resnet50_in_frozen/indist\endcsname{70.7}
\expandafter\def\csname valstore/yearbook/model/resnet50_in_frozen/future\endcsname{61.9}
\expandafter\def\csname valstore/yearbook/model/resnet50_in_frozen/decay\endcsname{8.7}
\expandafter\def\csname valstore/yearbook/model/resnet_l/indist\endcsname{90.0}
\expandafter\def\csname valstore/yearbook/model/resnet_l/future\endcsname{76.1}
\expandafter\def\csname valstore/yearbook/model/resnet_l/decay\endcsname{13.9}
\expandafter\def\csname valstore/yearbook/model/resnet_m/indist\endcsname{88.9}
\expandafter\def\csname valstore/yearbook/model/resnet_m/future\endcsname{75.8}
\expandafter\def\csname valstore/yearbook/model/resnet_m/decay\endcsname{13.1}
\expandafter\def\csname valstore/yearbook/model/resnet_s/indist\endcsname{87.3}
\expandafter\def\csname valstore/yearbook/model/resnet_s/future\endcsname{75.0}
\expandafter\def\csname valstore/yearbook/model/resnet_s/decay\endcsname{12.3}
\expandafter\def\csname valstore/yearbook/model/siglip_b_frozen/indist\endcsname{71.6}
\expandafter\def\csname valstore/yearbook/model/siglip_b_frozen/future\endcsname{62.4}
\expandafter\def\csname valstore/yearbook/model/siglip_b_frozen/decay\endcsname{9.2}
\expandafter\def\csname valstore/yearbook/model/vit_l/indist\endcsname{68.3}
\expandafter\def\csname valstore/yearbook/model/vit_l/future\endcsname{60.3}
\expandafter\def\csname valstore/yearbook/model/vit_l/decay\endcsname{8.0}
\expandafter\def\csname valstore/yearbook/model/vit_m/indist\endcsname{76.3}
\expandafter\def\csname valstore/yearbook/model/vit_m/future\endcsname{66.1}
\expandafter\def\csname valstore/yearbook/model/vit_m/decay\endcsname{10.1}
\expandafter\def\csname valstore/yearbook/model/vit_s/indist\endcsname{74.2}
\expandafter\def\csname valstore/yearbook/model/vit_s/future\endcsname{64.6}
\expandafter\def\csname valstore/yearbook/model/vit_s/decay\endcsname{9.6}
\expandafter\def\csname valstore/yearbook/model/vit_s16_in21k_frozen/indist\endcsname{72.2}
\expandafter\def\csname valstore/yearbook/model/vit_s16_in21k_frozen/future\endcsname{61.8}
\expandafter\def\csname valstore/yearbook/model/vit_s16_in21k_frozen/decay\endcsname{10.4}
\expandafter\def\csname valstore/yearbook/scratch/indist-min\endcsname{68.3}
\expandafter\def\csname valstore/yearbook/scratch/indist-max\endcsname{92.8}
\expandafter\def\csname valstore/yearbook/scratch/future-min\endcsname{60.3}
\expandafter\def\csname valstore/yearbook/scratch/future-max\endcsname{79.0}
\expandafter\def\csname valstore/yearbook/scratch/decay-min\endcsname{7.1}
\expandafter\def\csname valstore/yearbook/scratch/decay-max\endcsname{13.9}
\expandafter\def\csname valstore/yearbook/frozen/indist-min\endcsname{65.8}
\expandafter\def\csname valstore/yearbook/frozen/indist-max\endcsname{79.9}
\expandafter\def\csname valstore/yearbook/frozen/future-min\endcsname{57.9}
\expandafter\def\csname valstore/yearbook/frozen/future-max\endcsname{70.6}
\expandafter\def\csname valstore/yearbook/frozen/decay-min\endcsname{7.9}
\expandafter\def\csname valstore/yearbook/frozen/decay-max\endcsname{10.4}
\expandafter\def\csname valstore/yearbook/family/cnn/1905/future\endcsname{45.0}
\expandafter\def\csname valstore/yearbook/family/cnn/1905/decay\endcsname{8.3}
\expandafter\def\csname valstore/yearbook/family/cnn/1944/future\endcsname{80.3}
\expandafter\def\csname valstore/yearbook/family/cnn/1944/decay\endcsname{16.7}
\expandafter\def\csname valstore/yearbook/family/cnn/1978/future\endcsname{87.1}
\expandafter\def\csname valstore/yearbook/family/cnn/1978/decay\endcsname{-5.3}
\expandafter\def\csname valstore/yearbook/family/cnn/2012/future\endcsname{96.4}
\expandafter\def\csname valstore/yearbook/family/cnn/2012/decay\endcsname{-6.4}
\expandafter\def\csname valstore/yearbook/family/mlp/1905/future\endcsname{47.0}
\expandafter\def\csname valstore/yearbook/family/mlp/1905/decay\endcsname{-3.7}
\expandafter\def\csname valstore/yearbook/family/mlp/1944/future\endcsname{70.0}
\expandafter\def\csname valstore/yearbook/family/mlp/1944/decay\endcsname{12.4}
\expandafter\def\csname valstore/yearbook/family/mlp/1978/future\endcsname{79.0}
\expandafter\def\csname valstore/yearbook/family/mlp/1978/decay\endcsname{-4.8}
\expandafter\def\csname valstore/yearbook/family/mlp/2012/future\endcsname{86.6}
\expandafter\def\csname valstore/yearbook/family/mlp/2012/decay\endcsname{-5.3}
\expandafter\def\csname valstore/yearbook/family/resnet/1905/future\endcsname{49.9}
\expandafter\def\csname valstore/yearbook/family/resnet/1905/decay\endcsname{16.8}
\expandafter\def\csname valstore/yearbook/family/resnet/1944/future\endcsname{82.0}
\expandafter\def\csname valstore/yearbook/family/resnet/1944/decay\endcsname{16.2}
\expandafter\def\csname valstore/yearbook/family/resnet/1978/future\endcsname{88.0}
\expandafter\def\csname valstore/yearbook/family/resnet/1978/decay\endcsname{-1.4}
\expandafter\def\csname valstore/yearbook/family/resnet/2012/future\endcsname{96.6}
\expandafter\def\csname valstore/yearbook/family/resnet/2012/decay\endcsname{-7.1}
\expandafter\def\csname valstore/yearbook/family/transfer/1905/future\endcsname{46.9}
\expandafter\def\csname valstore/yearbook/family/transfer/1905/decay\endcsname{1.9}
\expandafter\def\csname valstore/yearbook/family/transfer/1944/future\endcsname{67.2}
\expandafter\def\csname valstore/yearbook/family/transfer/1944/decay\endcsname{21.4}
\expandafter\def\csname valstore/yearbook/family/transfer/1978/future\endcsname{70.6}
\expandafter\def\csname valstore/yearbook/family/transfer/1978/decay\endcsname{-4.4}
\expandafter\def\csname valstore/yearbook/family/transfer/2012/future\endcsname{85.4}
\expandafter\def\csname valstore/yearbook/family/transfer/2012/decay\endcsname{-6.1}
\expandafter\def\csname valstore/yearbook/family/vit/1905/future\endcsname{45.9}
\expandafter\def\csname valstore/yearbook/family/vit/1905/decay\endcsname{4.1}
\expandafter\def\csname valstore/yearbook/family/vit/1944/future\endcsname{70.7}
\expandafter\def\csname valstore/yearbook/family/vit/1944/decay\endcsname{20.3}
\expandafter\def\csname valstore/yearbook/family/vit/1978/future\endcsname{73.7}
\expandafter\def\csname valstore/yearbook/family/vit/1978/decay\endcsname{-5.9}
\expandafter\def\csname valstore/yearbook/family/vit/2012/future\endcsname{83.5}
\expandafter\def\csname valstore/yearbook/family/vit/2012/decay\endcsname{-8.4}
\expandafter\def\csname valstore/yearbook/family/transfer/indist\endcsname{71.8}
\expandafter\def\csname valstore/yearbook/family/transfer/future\endcsname{62.7}
\expandafter\def\csname valstore/yearbook/family/transfer/decay\endcsname{9.1}
\expandafter\def\csname valstore/yearbook/family/cnn/indist\endcsname{90.7}
\expandafter\def\csname valstore/yearbook/family/cnn/future\endcsname{77.4}
\expandafter\def\csname valstore/yearbook/family/cnn/decay\endcsname{13.3}
\expandafter\def\csname valstore/yearbook/family/mlp/indist\endcsname{80.0}
\expandafter\def\csname valstore/yearbook/family/mlp/future\endcsname{69.3}
\expandafter\def\csname valstore/yearbook/family/mlp/decay\endcsname{10.7}
\expandafter\def\csname valstore/yearbook/family/resnet/indist\endcsname{88.7}
\expandafter\def\csname valstore/yearbook/family/resnet/future\endcsname{75.6}
\expandafter\def\csname valstore/yearbook/family/resnet/decay\endcsname{13.1}
\expandafter\def\csname valstore/yearbook/family/vit/indist\endcsname{72.9}
\expandafter\def\csname valstore/yearbook/family/vit/future\endcsname{63.7}
\expandafter\def\csname valstore/yearbook/family/vit/decay\endcsname{9.2}
\expandafter\def\csname valstore/yearbook/param/mlp_s/trainable\endcsname{98k}
\expandafter\def\csname valstore/yearbook/param/mlp_s/total\endcsname{98k}
\expandafter\def\csname valstore/yearbook/param/mlp_m/trainable\endcsname{410k}
\expandafter\def\csname valstore/yearbook/param/mlp_m/total\endcsname{410k}
\expandafter\def\csname valstore/yearbook/param/mlp_l/trainable\endcsname{2.1M}
\expandafter\def\csname valstore/yearbook/param/mlp_l/total\endcsname{2.1M}
\expandafter\def\csname valstore/yearbook/param/cnn_s/trainable\endcsname{94k}
\expandafter\def\csname valstore/yearbook/param/cnn_s/total\endcsname{94k}
\expandafter\def\csname valstore/yearbook/param/cnn_m/trainable\endcsname{542k}
\expandafter\def\csname valstore/yearbook/param/cnn_m/total\endcsname{542k}
\expandafter\def\csname valstore/yearbook/param/cnn_l/trainable\endcsname{2.1M}
\expandafter\def\csname valstore/yearbook/param/cnn_l/total\endcsname{2.1M}
\expandafter\def\csname valstore/yearbook/param/resnet_s/trainable\endcsname{98k}
\expandafter\def\csname valstore/yearbook/param/resnet_s/total\endcsname{98k}
\expandafter\def\csname valstore/yearbook/param/resnet_m/trainable\endcsname{400k}
\expandafter\def\csname valstore/yearbook/param/resnet_m/total\endcsname{400k}
\expandafter\def\csname valstore/yearbook/param/resnet_l/trainable\endcsname{2.1M}
\expandafter\def\csname valstore/yearbook/param/resnet_l/total\endcsname{2.1M}
\expandafter\def\csname valstore/yearbook/param/vit_s/trainable\endcsname{75k}
\expandafter\def\csname valstore/yearbook/param/vit_s/total\endcsname{75k}
\expandafter\def\csname valstore/yearbook/param/vit_m/trainable\endcsname{545k}
\expandafter\def\csname valstore/yearbook/param/vit_m/total\endcsname{545k}
\expandafter\def\csname valstore/yearbook/param/vit_l/trainable\endcsname{2.2M}
\expandafter\def\csname valstore/yearbook/param/vit_l/total\endcsname{2.2M}
\expandafter\def\csname valstore/yearbook/param/resnet50_in_frozen/trainable\endcsname{4k}
\expandafter\def\csname valstore/yearbook/param/resnet50_in_frozen/total\endcsname{23.5M}
\expandafter\def\csname valstore/yearbook/param/vit_s16_in21k_frozen/trainable\endcsname{770}
\expandafter\def\csname valstore/yearbook/param/vit_s16_in21k_frozen/total\endcsname{21.7M}
\expandafter\def\csname valstore/yearbook/param/dinov2_s_frozen/trainable\endcsname{770}
\expandafter\def\csname valstore/yearbook/param/dinov2_s_frozen/total\endcsname{22.1M}
\expandafter\def\csname valstore/yearbook/param/dinov3_s_frozen/trainable\endcsname{770}
\expandafter\def\csname valstore/yearbook/param/dinov3_s_frozen/total\endcsname{21.6M}
\expandafter\def\csname valstore/yearbook/param/convnext_s_frozen/trainable\endcsname{2k}
\expandafter\def\csname valstore/yearbook/param/convnext_s_frozen/total\endcsname{49.5M}
\expandafter\def\csname valstore/yearbook/param/mae_b_frozen/trainable\endcsname{2k}
\expandafter\def\csname valstore/yearbook/param/mae_b_frozen/total\endcsname{85.8M}
\expandafter\def\csname valstore/yearbook/param/eva02_b_frozen/trainable\endcsname{2k}
\expandafter\def\csname valstore/yearbook/param/eva02_b_frozen/total\endcsname{85.8M}
\expandafter\def\csname valstore/yearbook/param/clip_b32_frozen/trainable\endcsname{2k}
\expandafter\def\csname valstore/yearbook/param/clip_b32_frozen/total\endcsname{87.5M}
\expandafter\def\csname valstore/yearbook/param/siglip_b_frozen/trainable\endcsname{2k}
\expandafter\def\csname valstore/yearbook/param/siglip_b_frozen/total\endcsname{92.9M}
% ---- Amazon Reviews ----
\expandafter\def\csname valstore/amazon_reviews/models\endcsname{19}
\expandafter\def\csname valstore/amazon_reviews/cutoffs\endcsname{20}
\expandafter\def\csname valstore/amazon_reviews/seeds\endcsname{3}
\expandafter\def\csname valstore/amazon_reviews/cutoff-first\endcsname{2014-H1}
\expandafter\def\csname valstore/amazon_reviews/cutoff-last\endcsname{2023-H2}
\expandafter\def\csname valstore/amazon_reviews/scratch-models\endcsname{12}
\expandafter\def\csname valstore/amazon_reviews/frozen-models\endcsname{7}
\expandafter\def\csname valstore/amazon_reviews/model/bert_base_frozen/indist\endcsname{0.924}
\expandafter\def\csname valstore/amazon_reviews/model/bert_base_frozen/future\endcsname{0.999}
\expandafter\def\csname valstore/amazon_reviews/model/bert_base_frozen/decay\endcsname{0.075}
\expandafter\def\csname valstore/amazon_reviews/model/bigru_s/indist\endcsname{0.680}
\expandafter\def\csname valstore/amazon_reviews/model/bigru_s/future\endcsname{0.797}
\expandafter\def\csname valstore/amazon_reviews/model/bigru_s/decay\endcsname{0.116}
\expandafter\def\csname valstore/amazon_reviews/model/bilstm_attn_l/indist\endcsname{0.713}
\expandafter\def\csname valstore/amazon_reviews/model/bilstm_attn_l/future\endcsname{0.819}
\expandafter\def\csname valstore/amazon_reviews/model/bilstm_attn_l/decay\endcsname{0.106}
\expandafter\def\csname valstore/amazon_reviews/model/bilstm_m/indist\endcsname{0.696}
\expandafter\def\csname valstore/amazon_reviews/model/bilstm_m/future\endcsname{0.824}
\expandafter\def\csname valstore/amazon_reviews/model/bilstm_m/decay\endcsname{0.128}
\expandafter\def\csname valstore/amazon_reviews/model/deberta_v3_base_frozen/indist\endcsname{0.806}
\expandafter\def\csname valstore/amazon_reviews/model/deberta_v3_base_frozen/future\endcsname{0.849}
\expandafter\def\csname valstore/amazon_reviews/model/deberta_v3_base_frozen/decay\endcsname{0.043}
\expandafter\def\csname valstore/amazon_reviews/model/distilbert_base_frozen/indist\endcsname{0.965}
\expandafter\def\csname valstore/amazon_reviews/model/distilbert_base_frozen/future\endcsname{1.047}
\expandafter\def\csname valstore/amazon_reviews/model/distilbert_base_frozen/decay\endcsname{0.082}
\expandafter\def\csname valstore/amazon_reviews/model/electra_base_frozen/indist\endcsname{0.924}
\expandafter\def\csname valstore/amazon_reviews/model/electra_base_frozen/future\endcsname{1.000}
\expandafter\def\csname valstore/amazon_reviews/model/electra_base_frozen/decay\endcsname{0.076}
\expandafter\def\csname valstore/amazon_reviews/model/ffn_l/indist\endcsname{0.759}
\expandafter\def\csname valstore/amazon_reviews/model/ffn_l/future\endcsname{0.848}
\expandafter\def\csname valstore/amazon_reviews/model/ffn_l/decay\endcsname{0.089}
\expandafter\def\csname valstore/amazon_reviews/model/ffn_m/indist\endcsname{0.757}
\expandafter\def\csname valstore/amazon_reviews/model/ffn_m/future\endcsname{0.834}
\expandafter\def\csname valstore/amazon_reviews/model/ffn_m/decay\endcsname{0.077}
\expandafter\def\csname valstore/amazon_reviews/model/ffn_s/indist\endcsname{0.766}
\expandafter\def\csname valstore/amazon_reviews/model/ffn_s/future\endcsname{0.841}
\expandafter\def\csname valstore/amazon_reviews/model/ffn_s/decay\endcsname{0.075}
\expandafter\def\csname valstore/amazon_reviews/model/modernbert_base_frozen/indist\endcsname{1.022}
\expandafter\def\csname valstore/amazon_reviews/model/modernbert_base_frozen/future\endcsname{1.086}
\expandafter\def\csname valstore/amazon_reviews/model/modernbert_base_frozen/decay\endcsname{0.064}
\expandafter\def\csname valstore/amazon_reviews/model/mpnet_base_frozen/indist\endcsname{0.948}
\expandafter\def\csname valstore/amazon_reviews/model/mpnet_base_frozen/future\endcsname{1.047}
\expandafter\def\csname valstore/amazon_reviews/model/mpnet_base_frozen/decay\endcsname{0.099}
\expandafter\def\csname valstore/amazon_reviews/model/roberta_base_frozen/indist\endcsname{0.856}
\expandafter\def\csname valstore/amazon_reviews/model/roberta_base_frozen/future\endcsname{0.944}
\expandafter\def\csname valstore/amazon_reviews/model/roberta_base_frozen/decay\endcsname{0.088}
\expandafter\def\csname valstore/amazon_reviews/model/textcnn_l/indist\endcsname{0.793}
\expandafter\def\csname valstore/amazon_reviews/model/textcnn_l/future\endcsname{0.871}
\expandafter\def\csname valstore/amazon_reviews/model/textcnn_l/decay\endcsname{0.078}
\expandafter\def\csname valstore/amazon_reviews/model/textcnn_m/indist\endcsname{0.792}
\expandafter\def\csname valstore/amazon_reviews/model/textcnn_m/future\endcsname{0.885}
\expandafter\def\csname valstore/amazon_reviews/model/textcnn_m/decay\endcsname{0.093}
\expandafter\def\csname valstore/amazon_reviews/model/textcnn_s/indist\endcsname{0.746}
\expandafter\def\csname valstore/amazon_reviews/model/textcnn_s/future\endcsname{0.860}
\expandafter\def\csname valstore/amazon_reviews/model/textcnn_s/decay\endcsname{0.114}
\expandafter\def\csname valstore/amazon_reviews/model/tx_l/indist\endcsname{0.735}
\expandafter\def\csname valstore/amazon_reviews/model/tx_l/future\endcsname{0.829}
\expandafter\def\csname valstore/amazon_reviews/model/tx_l/decay\endcsname{0.094}
\expandafter\def\csname valstore/amazon_reviews/model/tx_m/indist\endcsname{0.708}
\expandafter\def\csname valstore/amazon_reviews/model/tx_m/future\endcsname{0.801}
\expandafter\def\csname valstore/amazon_reviews/model/tx_m/decay\endcsname{0.092}
\expandafter\def\csname valstore/amazon_reviews/model/tx_s/indist\endcsname{0.686}
\expandafter\def\csname valstore/amazon_reviews/model/tx_s/future\endcsname{0.776}
\expandafter\def\csname valstore/amazon_reviews/model/tx_s/decay\endcsname{0.090}
\expandafter\def\csname valstore/amazon_reviews/scratch/indist-min\endcsname{0.680}
\expandafter\def\csname valstore/amazon_reviews/scratch/indist-max\endcsname{0.793}
\expandafter\def\csname valstore/amazon_reviews/scratch/future-min\endcsname{0.776}
\expandafter\def\csname valstore/amazon_reviews/scratch/future-max\endcsname{0.885}
\expandafter\def\csname valstore/amazon_reviews/scratch/decay-min\endcsname{0.075}
\expandafter\def\csname valstore/amazon_reviews/scratch/decay-max\endcsname{0.128}
\expandafter\def\csname valstore/amazon_reviews/frozen/indist-min\endcsname{0.806}
\expandafter\def\csname valstore/amazon_reviews/frozen/indist-max\endcsname{1.022}
\expandafter\def\csname valstore/amazon_reviews/frozen/future-min\endcsname{0.849}
\expandafter\def\csname valstore/amazon_reviews/frozen/future-max\endcsname{1.086}
\expandafter\def\csname valstore/amazon_reviews/frozen/decay-min\endcsname{0.043}
\expandafter\def\csname valstore/amazon_reviews/frozen/decay-max\endcsname{0.099}
\expandafter\def\csname valstore/amazon_reviews/family/ffn/28/future\endcsname{0.993}
\expandafter\def\csname valstore/amazon_reviews/family/ffn/28/decay\endcsname{0.117}
\expandafter\def\csname valstore/amazon_reviews/family/ffn/34/future\endcsname{0.817}
\expandafter\def\csname valstore/amazon_reviews/family/ffn/34/decay\endcsname{0.087}
\expandafter\def\csname valstore/amazon_reviews/family/ffn/40/future\endcsname{0.802}
\expandafter\def\csname valstore/amazon_reviews/family/ffn/40/decay\endcsname{0.073}
\expandafter\def\csname valstore/amazon_reviews/family/ffn/46/future\endcsname{0.789}
\expandafter\def\csname valstore/amazon_reviews/family/ffn/46/decay\endcsname{0.056}
\expandafter\def\csname valstore/amazon_reviews/family/frozen/28/future\endcsname{1.351}
\expandafter\def\csname valstore/amazon_reviews/family/frozen/28/decay\endcsname{0.131}
\expandafter\def\csname valstore/amazon_reviews/family/frozen/34/future\endcsname{0.943}
\expandafter\def\csname valstore/amazon_reviews/family/frozen/34/decay\endcsname{0.068}
\expandafter\def\csname valstore/amazon_reviews/family/frozen/40/future\endcsname{0.934}
\expandafter\def\csname valstore/amazon_reviews/family/frozen/40/decay\endcsname{0.073}
\expandafter\def\csname valstore/amazon_reviews/family/frozen/46/future\endcsname{0.921}
\expandafter\def\csname valstore/amazon_reviews/family/frozen/46/decay\endcsname{0.039}
\expandafter\def\csname valstore/amazon_reviews/family/recurrent/28/future\endcsname{1.101}
\expandafter\def\csname valstore/amazon_reviews/family/recurrent/28/decay\endcsname{0.242}
\expandafter\def\csname valstore/amazon_reviews/family/recurrent/34/future\endcsname{0.794}
\expandafter\def\csname valstore/amazon_reviews/family/recurrent/34/decay\endcsname{0.134}
\expandafter\def\csname valstore/amazon_reviews/family/recurrent/40/future\endcsname{0.723}
\expandafter\def\csname valstore/amazon_reviews/family/recurrent/40/decay\endcsname{0.094}
\expandafter\def\csname valstore/amazon_reviews/family/recurrent/46/future\endcsname{0.710}
\expandafter\def\csname valstore/amazon_reviews/family/recurrent/46/decay\endcsname{0.059}
\expandafter\def\csname valstore/amazon_reviews/family/textcnn/28/future\endcsname{1.002}
\expandafter\def\csname valstore/amazon_reviews/family/textcnn/28/decay\endcsname{0.227}
\expandafter\def\csname valstore/amazon_reviews/family/textcnn/34/future\endcsname{0.859}
\expandafter\def\csname valstore/amazon_reviews/family/textcnn/34/decay\endcsname{0.102}
\expandafter\def\csname valstore/amazon_reviews/family/textcnn/40/future\endcsname{0.811}
\expandafter\def\csname valstore/amazon_reviews/family/textcnn/40/decay\endcsname{0.090}
\expandafter\def\csname valstore/amazon_reviews/family/textcnn/46/future\endcsname{0.818}
\expandafter\def\csname valstore/amazon_reviews/family/textcnn/46/decay\endcsname{0.076}
\expandafter\def\csname valstore/amazon_reviews/family/transformer/28/future\endcsname{0.967}
\expandafter\def\csname valstore/amazon_reviews/family/transformer/28/decay\endcsname{0.173}
\expandafter\def\csname valstore/amazon_reviews/family/transformer/34/future\endcsname{0.784}
\expandafter\def\csname valstore/amazon_reviews/family/transformer/34/decay\endcsname{0.101}
\expandafter\def\csname valstore/amazon_reviews/family/transformer/40/future\endcsname{0.748}
\expandafter\def\csname valstore/amazon_reviews/family/transformer/40/decay\endcsname{0.080}
\expandafter\def\csname valstore/amazon_reviews/family/transformer/46/future\endcsname{0.721}
\expandafter\def\csname valstore/amazon_reviews/family/transformer/46/decay\endcsname{0.056}
\expandafter\def\csname valstore/amazon_reviews/family/frozen/indist\endcsname{0.921}
\expandafter\def\csname valstore/amazon_reviews/family/frozen/future\endcsname{0.996}
\expandafter\def\csname valstore/amazon_reviews/family/frozen/decay\endcsname{0.075}
\expandafter\def\csname valstore/amazon_reviews/family/recurrent/indist\endcsname{0.696}
\expandafter\def\csname valstore/amazon_reviews/family/recurrent/future\endcsname{0.813}
\expandafter\def\csname valstore/amazon_reviews/family/recurrent/decay\endcsname{0.117}
\expandafter\def\csname valstore/amazon_reviews/family/ffn/indist\endcsname{0.761}
\expandafter\def\csname valstore/amazon_reviews/family/ffn/future\endcsname{0.841}
\expandafter\def\csname valstore/amazon_reviews/family/ffn/decay\endcsname{0.080}
\expandafter\def\csname valstore/amazon_reviews/family/textcnn/indist\endcsname{0.777}
\expandafter\def\csname valstore/amazon_reviews/family/textcnn/future\endcsname{0.872}
\expandafter\def\csname valstore/amazon_reviews/family/textcnn/decay\endcsname{0.095}
\expandafter\def\csname valstore/amazon_reviews/family/transformer/indist\endcsname{0.710}
\expandafter\def\csname valstore/amazon_reviews/family/transformer/future\endcsname{0.802}
\expandafter\def\csname valstore/amazon_reviews/family/transformer/decay\endcsname{0.092}
\expandafter\def\csname valstore/amazon_reviews/param/ffn_s/trainable\endcsname{99k}
\expandafter\def\csname valstore/amazon_reviews/param/ffn_s/total\endcsname{124.7M}
\expandafter\def\csname valstore/amazon_reviews/param/ffn_m/trainable\endcsname{394k}
\expandafter\def\csname valstore/amazon_reviews/param/ffn_m/total\endcsname{125.0M}
\expandafter\def\csname valstore/amazon_reviews/param/ffn_l/trainable\endcsname{1.6M}
\expandafter\def\csname valstore/amazon_reviews/param/ffn_l/total\endcsname{126.2M}
\expandafter\def\csname valstore/amazon_reviews/param/textcnn_s/trainable\endcsname{88k}
\expandafter\def\csname valstore/amazon_reviews/param/textcnn_s/total\endcsname{124.7M}
\expandafter\def\csname valstore/amazon_reviews/param/textcnn_m/trainable\endcsname{461k}
\expandafter\def\csname valstore/amazon_reviews/param/textcnn_m/total\endcsname{125.1M}
\expandafter\def\csname valstore/amazon_reviews/param/textcnn_l/trainable\endcsname{1.9M}
\expandafter\def\csname valstore/amazon_reviews/param/textcnn_l/total\endcsname{126.6M}
\expandafter\def\csname valstore/amazon_reviews/param/bigru_s/trainable\endcsname{99k}
\expandafter\def\csname valstore/amazon_reviews/param/bigru_s/total\endcsname{124.7M}
\expandafter\def\csname valstore/amazon_reviews/param/bilstm_m/trainable\endcsname{536k}
\expandafter\def\csname valstore/amazon_reviews/param/bilstm_m/total\endcsname{125.2M}
\expandafter\def\csname valstore/amazon_reviews/param/bilstm_attn_l/trainable\endcsname{2.2M}
\expandafter\def\csname valstore/amazon_reviews/param/bilstm_attn_l/total\endcsname{126.8M}
\expandafter\def\csname valstore/amazon_reviews/param/tx_s/trainable\endcsname{83k}
\expandafter\def\csname valstore/amazon_reviews/param/tx_s/total\endcsname{124.7M}
\expandafter\def\csname valstore/amazon_reviews/param/tx_m/trainable\endcsname{493k}
\expandafter\def\csname valstore/amazon_reviews/param/tx_m/total\endcsname{125.1M}
\expandafter\def\csname valstore/amazon_reviews/param/tx_l/trainable\endcsname{1.9M}
\expandafter\def\csname valstore/amazon_reviews/param/tx_l/total\endcsname{126.5M}
\expandafter\def\csname valstore/amazon_reviews/param/distilbert_base_frozen/trainable\endcsname{769}
\expandafter\def\csname valstore/amazon_reviews/param/distilbert_base_frozen/total\endcsname{66.4M}
\expandafter\def\csname valstore/amazon_reviews/param/bert_base_frozen/trainable\endcsname{769}
\expandafter\def\csname valstore/amazon_reviews/param/bert_base_frozen/total\endcsname{109.5M}
\expandafter\def\csname valstore/amazon_reviews/param/roberta_base_frozen/trainable\endcsname{769}
\expandafter\def\csname valstore/amazon_reviews/param/roberta_base_frozen/total\endcsname{124.6M}
\expandafter\def\csname valstore/amazon_reviews/param/deberta_v3_base_frozen/trainable\endcsname{769}
\expandafter\def\csname valstore/amazon_reviews/param/deberta_v3_base_frozen/total\endcsname{183.8M}
\expandafter\def\csname valstore/amazon_reviews/param/electra_base_frozen/trainable\endcsname{769}
\expandafter\def\csname valstore/amazon_reviews/param/electra_base_frozen/total\endcsname{108.9M}
\expandafter\def\csname valstore/amazon_reviews/param/mpnet_base_frozen/trainable\endcsname{769}
\expandafter\def\csname valstore/amazon_reviews/param/mpnet_base_frozen/total\endcsname{109.5M}
\expandafter\def\csname valstore/amazon_reviews/param/modernbert_base_frozen/trainable\endcsname{769}
\expandafter\def\csname valstore/amazon_reviews/param/modernbert_base_frozen/total\endcsname{149.0M}
% ---- arXiv ----
\expandafter\def\csname valstore/arxiv/models\endcsname{20}
\expandafter\def\csname valstore/arxiv/cutoffs\endcsname{26}
\expandafter\def\csname valstore/arxiv/seeds\endcsname{3}
\expandafter\def\csname valstore/arxiv/cutoff-first\endcsname{2000}
\expandafter\def\csname valstore/arxiv/cutoff-last\endcsname{2025}
\expandafter\def\csname valstore/arxiv/scratch-models\endcsname{12}
\expandafter\def\csname valstore/arxiv/frozen-models\endcsname{8}
\expandafter\def\csname valstore/arxiv/model/bert_base_frozen/indist\endcsname{96.7}
\expandafter\def\csname valstore/arxiv/model/bert_base_frozen/future\endcsname{93.6}
\expandafter\def\csname valstore/arxiv/model/bert_base_frozen/decay\endcsname{3.1}
\expandafter\def\csname valstore/arxiv/model/bigru_s/indist\endcsname{97.9}
\expandafter\def\csname valstore/arxiv/model/bigru_s/future\endcsname{94.7}
\expandafter\def\csname valstore/arxiv/model/bigru_s/decay\endcsname{3.2}
\expandafter\def\csname valstore/arxiv/model/bilstm_attn_l/indist\endcsname{98.0}
\expandafter\def\csname valstore/arxiv/model/bilstm_attn_l/future\endcsname{95.1}
\expandafter\def\csname valstore/arxiv/model/bilstm_attn_l/decay\endcsname{2.9}
\expandafter\def\csname valstore/arxiv/model/bilstm_m/indist\endcsname{98.1}
\expandafter\def\csname valstore/arxiv/model/bilstm_m/future\endcsname{95.2}
\expandafter\def\csname valstore/arxiv/model/bilstm_m/decay\endcsname{2.9}
\expandafter\def\csname valstore/arxiv/model/deberta_v3_base_frozen/indist\endcsname{92.7}
\expandafter\def\csname valstore/arxiv/model/deberta_v3_base_frozen/future\endcsname{86.1}
\expandafter\def\csname valstore/arxiv/model/deberta_v3_base_frozen/decay\endcsname{6.5}
\expandafter\def\csname valstore/arxiv/model/distilbert_base_frozen/indist\endcsname{96.4}
\expandafter\def\csname valstore/arxiv/model/distilbert_base_frozen/future\endcsname{93.1}
\expandafter\def\csname valstore/arxiv/model/distilbert_base_frozen/decay\endcsname{3.3}
\expandafter\def\csname valstore/arxiv/model/electra_base_frozen/indist\endcsname{91.8}
\expandafter\def\csname valstore/arxiv/model/electra_base_frozen/future\endcsname{84.4}
\expandafter\def\csname valstore/arxiv/model/electra_base_frozen/decay\endcsname{7.3}
\expandafter\def\csname valstore/arxiv/model/ffn_l/indist\endcsname{97.6}
\expandafter\def\csname valstore/arxiv/model/ffn_l/future\endcsname{94.9}
\expandafter\def\csname valstore/arxiv/model/ffn_l/decay\endcsname{2.7}
\expandafter\def\csname valstore/arxiv/model/ffn_m/indist\endcsname{97.5}
\expandafter\def\csname valstore/arxiv/model/ffn_m/future\endcsname{94.6}
\expandafter\def\csname valstore/arxiv/model/ffn_m/decay\endcsname{2.9}
\expandafter\def\csname valstore/arxiv/model/ffn_s/indist\endcsname{97.2}
\expandafter\def\csname valstore/arxiv/model/ffn_s/future\endcsname{93.9}
\expandafter\def\csname valstore/arxiv/model/ffn_s/decay\endcsname{3.4}
\expandafter\def\csname valstore/arxiv/model/minilm_l6_frozen/indist\endcsname{97.7}
\expandafter\def\csname valstore/arxiv/model/minilm_l6_frozen/future\endcsname{94.8}
\expandafter\def\csname valstore/arxiv/model/minilm_l6_frozen/decay\endcsname{2.8}
\expandafter\def\csname valstore/arxiv/model/modernbert_base_frozen/indist\endcsname{97.4}
\expandafter\def\csname valstore/arxiv/model/modernbert_base_frozen/future\endcsname{94.9}
\expandafter\def\csname valstore/arxiv/model/modernbert_base_frozen/decay\endcsname{2.5}
\expandafter\def\csname valstore/arxiv/model/mpnet_base_frozen/indist\endcsname{95.7}
\expandafter\def\csname valstore/arxiv/model/mpnet_base_frozen/future\endcsname{92.0}
\expandafter\def\csname valstore/arxiv/model/mpnet_base_frozen/decay\endcsname{3.7}
\expandafter\def\csname valstore/arxiv/model/roberta_base_frozen/indist\endcsname{95.9}
\expandafter\def\csname valstore/arxiv/model/roberta_base_frozen/future\endcsname{92.3}
\expandafter\def\csname valstore/arxiv/model/roberta_base_frozen/decay\endcsname{3.6}
\expandafter\def\csname valstore/arxiv/model/textcnn_l/indist\endcsname{98.0}
\expandafter\def\csname valstore/arxiv/model/textcnn_l/future\endcsname{95.2}
\expandafter\def\csname valstore/arxiv/model/textcnn_l/decay\endcsname{2.7}
\expandafter\def\csname valstore/arxiv/model/textcnn_m/indist\endcsname{97.9}
\expandafter\def\csname valstore/arxiv/model/textcnn_m/future\endcsname{94.8}
\expandafter\def\csname valstore/arxiv/model/textcnn_m/decay\endcsname{3.1}
\expandafter\def\csname valstore/arxiv/model/textcnn_s/indist\endcsname{97.9}
\expandafter\def\csname valstore/arxiv/model/textcnn_s/future\endcsname{94.5}
\expandafter\def\csname valstore/arxiv/model/textcnn_s/decay\endcsname{3.4}
\expandafter\def\csname valstore/arxiv/model/tx_l/indist\endcsname{97.9}
\expandafter\def\csname valstore/arxiv/model/tx_l/future\endcsname{94.7}
\expandafter\def\csname valstore/arxiv/model/tx_l/decay\endcsname{3.1}
\expandafter\def\csname valstore/arxiv/model/tx_m/indist\endcsname{97.9}
\expandafter\def\csname valstore/arxiv/model/tx_m/future\endcsname{95.1}
\expandafter\def\csname valstore/arxiv/model/tx_m/decay\endcsname{2.8}
\expandafter\def\csname valstore/arxiv/model/tx_s/indist\endcsname{98.0}
\expandafter\def\csname valstore/arxiv/model/tx_s/future\endcsname{95.1}
\expandafter\def\csname valstore/arxiv/model/tx_s/decay\endcsname{2.9}
\expandafter\def\csname valstore/arxiv/scratch/indist-min\endcsname{97.2}
\expandafter\def\csname valstore/arxiv/scratch/indist-max\endcsname{98.1}
\expandafter\def\csname valstore/arxiv/scratch/future-min\endcsname{93.9}
\expandafter\def\csname valstore/arxiv/scratch/future-max\endcsname{95.2}
\expandafter\def\csname valstore/arxiv/scratch/decay-min\endcsname{2.7}
\expandafter\def\csname valstore/arxiv/scratch/decay-max\endcsname{3.4}
\expandafter\def\csname valstore/arxiv/frozen/indist-min\endcsname{91.8}
\expandafter\def\csname valstore/arxiv/frozen/indist-max\endcsname{97.7}
\expandafter\def\csname valstore/arxiv/frozen/future-min\endcsname{84.4}
\expandafter\def\csname valstore/arxiv/frozen/future-max\endcsname{94.9}
\expandafter\def\csname valstore/arxiv/frozen/decay-min\endcsname{2.5}
\expandafter\def\csname valstore/arxiv/frozen/decay-max\endcsname{7.3}
\expandafter\def\csname valstore/arxiv/family/ffn/2000/future\endcsname{92.2}
\expandafter\def\csname valstore/arxiv/family/ffn/2000/decay\endcsname{4.3}
\expandafter\def\csname valstore/arxiv/family/ffn/2008/future\endcsname{94.6}
\expandafter\def\csname valstore/arxiv/family/ffn/2008/decay\endcsname{3.3}
\expandafter\def\csname valstore/arxiv/family/ffn/2016/future\endcsname{96.1}
\expandafter\def\csname valstore/arxiv/family/ffn/2016/decay\endcsname{2.2}
\expandafter\def\csname valstore/arxiv/family/ffn/2024/future\endcsname{95.8}
\expandafter\def\csname valstore/arxiv/family/ffn/2024/decay\endcsname{0.2}
\expandafter\def\csname valstore/arxiv/family/frozen/2000/future\endcsname{87.8}
\expandafter\def\csname valstore/arxiv/family/frozen/2000/decay\endcsname{5.4}
\expandafter\def\csname valstore/arxiv/family/frozen/2008/future\endcsname{91.6}
\expandafter\def\csname valstore/arxiv/family/frozen/2008/decay\endcsname{4.6}
\expandafter\def\csname valstore/arxiv/family/frozen/2016/future\endcsname{93.6}
\expandafter\def\csname valstore/arxiv/family/frozen/2016/decay\endcsname{2.8}
\expandafter\def\csname valstore/arxiv/family/frozen/2024/future\endcsname{93.8}
\expandafter\def\csname valstore/arxiv/family/frozen/2024/decay\endcsname{0.5}
\expandafter\def\csname valstore/arxiv/family/recurrent/2000/future\endcsname{92.7}
\expandafter\def\csname valstore/arxiv/family/recurrent/2000/decay\endcsname{3.4}
\expandafter\def\csname valstore/arxiv/family/recurrent/2008/future\endcsname{95.2}
\expandafter\def\csname valstore/arxiv/family/recurrent/2008/decay\endcsname{3.4}
\expandafter\def\csname valstore/arxiv/family/recurrent/2016/future\endcsname{96.5}
\expandafter\def\csname valstore/arxiv/family/recurrent/2016/decay\endcsname{2.2}
\expandafter\def\csname valstore/arxiv/family/recurrent/2024/future\endcsname{96.3}
\expandafter\def\csname valstore/arxiv/family/recurrent/2024/decay\endcsname{0.2}
\expandafter\def\csname valstore/arxiv/family/textcnn/2000/future\endcsname{93.4}
\expandafter\def\csname valstore/arxiv/family/textcnn/2000/decay\endcsname{4.3}
\expandafter\def\csname valstore/arxiv/family/textcnn/2008/future\endcsname{94.7}
\expandafter\def\csname valstore/arxiv/family/textcnn/2008/decay\endcsname{3.9}
\expandafter\def\csname valstore/arxiv/family/textcnn/2016/future\endcsname{96.2}
\expandafter\def\csname valstore/arxiv/family/textcnn/2016/decay\endcsname{2.3}
\expandafter\def\csname valstore/arxiv/family/textcnn/2024/future\endcsname{95.8}
\expandafter\def\csname valstore/arxiv/family/textcnn/2024/decay\endcsname{0.3}
\expandafter\def\csname valstore/arxiv/family/transformer/2000/future\endcsname{93.5}
\expandafter\def\csname valstore/arxiv/family/transformer/2000/decay\endcsname{2.9}
\expandafter\def\csname valstore/arxiv/family/transformer/2008/future\endcsname{95.4}
\expandafter\def\csname valstore/arxiv/family/transformer/2008/decay\endcsname{3.0}
\expandafter\def\csname valstore/arxiv/family/transformer/2016/future\endcsname{96.5}
\expandafter\def\csname valstore/arxiv/family/transformer/2016/decay\endcsname{2.2}
\expandafter\def\csname valstore/arxiv/family/transformer/2024/future\endcsname{96.3}
\expandafter\def\csname valstore/arxiv/family/transformer/2024/decay\endcsname{0.2}
\expandafter\def\csname valstore/arxiv/family/frozen/indist\endcsname{95.5}
\expandafter\def\csname valstore/arxiv/family/frozen/future\endcsname{91.4}
\expandafter\def\csname valstore/arxiv/family/frozen/decay\endcsname{4.1}
\expandafter\def\csname valstore/arxiv/family/recurrent/indist\endcsname{98.0}
\expandafter\def\csname valstore/arxiv/family/recurrent/future\endcsname{95.0}
\expandafter\def\csname valstore/arxiv/family/recurrent/decay\endcsname{3.0}
\expandafter\def\csname valstore/arxiv/family/ffn/indist\endcsname{97.4}
\expandafter\def\csname valstore/arxiv/family/ffn/future\endcsname{94.4}
\expandafter\def\csname valstore/arxiv/family/ffn/decay\endcsname{3.0}
\expandafter\def\csname valstore/arxiv/family/textcnn/indist\endcsname{97.9}
\expandafter\def\csname valstore/arxiv/family/textcnn/future\endcsname{94.8}
\expandafter\def\csname valstore/arxiv/family/textcnn/decay\endcsname{3.1}
\expandafter\def\csname valstore/arxiv/family/transformer/indist\endcsname{97.9}
\expandafter\def\csname valstore/arxiv/family/transformer/future\endcsname{95.0}
\expandafter\def\csname valstore/arxiv/family/transformer/decay\endcsname{2.9}
\expandafter\def\csname valstore/arxiv/param/ffn_s/trainable\endcsname{99k}
\expandafter\def\csname valstore/arxiv/param/ffn_s/total\endcsname{124.7M}
\expandafter\def\csname valstore/arxiv/param/ffn_m/trainable\endcsname{397k}
\expandafter\def\csname valstore/arxiv/param/ffn_m/total\endcsname{125.0M}
\expandafter\def\csname valstore/arxiv/param/ffn_l/trainable\endcsname{1.6M}
\expandafter\def\csname valstore/arxiv/param/ffn_l/total\endcsname{126.2M}
\expandafter\def\csname valstore/arxiv/param/textcnn_s/trainable\endcsname{89k}
\expandafter\def\csname valstore/arxiv/param/textcnn_s/total\endcsname{124.7M}
\expandafter\def\csname valstore/arxiv/param/textcnn_m/trainable\endcsname{462k}
\expandafter\def\csname valstore/arxiv/param/textcnn_m/total\endcsname{125.1M}
\expandafter\def\csname valstore/arxiv/param/textcnn_l/trainable\endcsname{1.9M}
\expandafter\def\csname valstore/arxiv/param/textcnn_l/total\endcsname{126.6M}
\expandafter\def\csname valstore/arxiv/param/bigru_s/trainable\endcsname{100k}
\expandafter\def\csname valstore/arxiv/param/bigru_s/total\endcsname{124.7M}
\expandafter\def\csname valstore/arxiv/param/bilstm_m/trainable\endcsname{537k}
\expandafter\def\csname valstore/arxiv/param/bilstm_m/total\endcsname{125.2M}
\expandafter\def\csname valstore/arxiv/param/bilstm_attn_l/trainable\endcsname{2.2M}
\expandafter\def\csname valstore/arxiv/param/bilstm_attn_l/total\endcsname{126.8M}
\expandafter\def\csname valstore/arxiv/param/tx_s/trainable\endcsname{83k}
\expandafter\def\csname valstore/arxiv/param/tx_s/total\endcsname{124.7M}
\expandafter\def\csname valstore/arxiv/param/tx_m/trainable\endcsname{493k}
\expandafter\def\csname valstore/arxiv/param/tx_m/total\endcsname{125.1M}
\expandafter\def\csname valstore/arxiv/param/tx_l/trainable\endcsname{1.9M}
\expandafter\def\csname valstore/arxiv/param/tx_l/total\endcsname{126.5M}
\expandafter\def\csname valstore/arxiv/param/minilm_l6_frozen/trainable\endcsname{3k}
\expandafter\def\csname valstore/arxiv/param/minilm_l6_frozen/total\endcsname{22.7M}
\expandafter\def\csname valstore/arxiv/param/distilbert_base_frozen/trainable\endcsname{5k}
\expandafter\def\csname valstore/arxiv/param/distilbert_base_frozen/total\endcsname{66.4M}
\expandafter\def\csname valstore/arxiv/param/bert_base_frozen/trainable\endcsname{5k}
\expandafter\def\csname valstore/arxiv/param/bert_base_frozen/total\endcsname{109.5M}
\expandafter\def\csname valstore/arxiv/param/roberta_base_frozen/trainable\endcsname{5k}
\expandafter\def\csname valstore/arxiv/param/roberta_base_frozen/total\endcsname{124.7M}
\expandafter\def\csname valstore/arxiv/param/deberta_v3_base_frozen/trainable\endcsname{5k}
\expandafter\def\csname valstore/arxiv/param/deberta_v3_base_frozen/total\endcsname{183.8M}
\expandafter\def\csname valstore/arxiv/param/electra_base_frozen/trainable\endcsname{5k}
\expandafter\def\csname valstore/arxiv/param/electra_base_frozen/total\endcsname{108.9M}
\expandafter\def\csname valstore/arxiv/param/mpnet_base_frozen/trainable\endcsname{5k}
\expandafter\def\csname valstore/arxiv/param/mpnet_base_frozen/total\endcsname{109.5M}
\expandafter\def\csname valstore/arxiv/param/modernbert_base_frozen/trainable\endcsname{5k}
\expandafter\def\csname valstore/arxiv/param/modernbert_base_frozen/total\endcsname{149.0M}


%%%%%%%%%%%%%%%%%%%%%%%%%%%%%%%%
% Figure panels
%%%%%%%%%%%%%%%%%%%%%%%%%%%%%%%%
\newenvironment{figurepanel}[1]{%
  \begin{minipage}[t]{#1}\centering
}{%
  \end{minipage}%
}
\newcommand{\panellabel}[1]{\par\vspace{0.25ex}{\footnotesize #1}\par}

%%%%%%%%%%%%%%%%%%%%%%%%%%%%%%%%
% Hyperlinks and clever references (loaded last, in this order)
%%%%%%%%%%%%%%%%%%%%%%%%%%%%%%%%
\usepackage{hyperref}
\usepackage[capitalize,noabbrev]{cleveref}
\definecolor{linknavy}{rgb}{0,0.08,0.45}
\hypersetup{colorlinks=true,linkcolor=linknavy,citecolor=linknavy,urlcolor=linknavy,pdfborder={0 0 0}}

%%%%%%%%%%%%%%%%%%%%%%%%%%%%%%%%
% Single-source hooks
%%%%%%%%%%%%%%%%%%%%%%%%%%%%%%%%
% Build discriminator: true only in the page-limited LNCS proceedings build.
% Lets shared source drop content that the roomy preprint keeps, e.g. extra
% citations and the duplicate Yearbook reference. Defaults to false here; the
% LNCS root sets \lncstrue after \input-ing this file.
\newif\iflncs

% Compact vertical spacing used only where a dense two-column layout needs
% it. A no-op by default (roomy layouts such as LNCS); the preprint root
% redefines it to \vspace so the shared content stays single-sourced.
\providecommand{\tightvspace}[1]{}

% Verbose display equations. The extended preprint spells out the standard
% loss and drift-summary formulas as numbered display equations; the
% page-constrained LNCS build inlines the same definitions. Default off
% (LNCS compact); the preprint root switches it on with \verbosemathtrue.
\newif\ifverbosemath
\verbosemathfalse

% Durable links used when the proceedings build points readers to the extended
% evidence bundle kept outside the compact LNCS appendix.
\newcommand{\extendedpreprintlink}{\href{https://drift-happens.org/drift-happens.pdf}{extended preprint}}
\newcommand{\resultssitelink}{\href{https://drift-happens.org}{drift-happens.org}}

% Keyword list shared by both builds; consumed by each root as its format
% requires (LNCS \keywords{}, preprint metadata / printed line).
\newcommand{\abskeywords}{Temporal distribution shift \and Concept drift \and Inductive bias \and Neural architecture robustness \and Pretrained encoders}

% Main-body forgetting curve. The roomy right-hand-legend plot reads well in
% the two-column / web layouts, so it is the default here; the single-column
% LNCS root overrides these to the wide, short legend-below variant that stays
% readable at text width and packs without a stranded half-page (see
% main-lncs.tex). The float opener carries the placement, since LaTeX does not
% expand macros inside a float's placement argument.
\providecommand{\forgettingfloat}{\begin{figure}[H]}
\providecommand{\forgettinggraphic}{%
  \includegraphics[width=0.45\textwidth]{plots_experiments/yearbook/forgetting.pdf}}
